\documentclass[11pt]{article}

% The document is written for a modern LaTeX engine.  The bibliography is
% intentionally kept external because the original references.bib was not
% included with the pasted report.
\usepackage[T1]{fontenc}
\usepackage[margin=1in]{geometry}
\usepackage{microtype}
\usepackage{amsmath,amssymb,mathtools}
\usepackage{booktabs}
\usepackage{graphicx}
\usepackage{csquotes}
\usepackage[backend=biber,style=numeric,sorting=none]{biblatex}
\usepackage[hidelinks]{hyperref}
\usepackage[nameinlink,noabbrev]{cleveref}

\addbibresource{references.bib}

% Reusable notation keeps the mathematical statements compact and consistent.
\newcommand{\R}{\mathbb{R}}
\newcommand{\latent}{\mathcal{Z}}
\newcommand{\structure}{\mathcal{X}}
\newcommand{\sphere}{\mathbb{S}}
\newcommand{\norm}[1]{\lVert #1 \rVert_2}
\newcommand{\TODO}[1]{\textit{[To be completed: #1]}}
\DeclareMathOperator{\TMScore}{TM\text{-}score}
\DeclareMathOperator{\LogEI}{LogEI}

\title{September 4, 2026 Report\\[0.5em]
       \large Surrogate Latent Spaces for Budgeted Black-Box Optimization}
\author{Thanh Tran}
\date{August 2026}

\begin{document}

\maketitle

\begin{abstract}
This report summarizes the construction of a low-dimensional surrogate latent
space for black-box optimization of a frozen generative model.  The central
idea is to represent candidate latents as unit-norm combinations of selected
seed latents, then use a measure-preserving Knothe--Rosenblatt chart to give
the resulting hyperspherical search space Euclidean coordinates.  The report
focuses on the mathematical construction and its implementation assumptions;
experimental outcomes are intentionally left as placeholders.
\end{abstract}

\section{Purpose and scope}

This report revisits the methods introduced by Dr.~Huang and studies their
application to sample-efficient black-box optimization of generative-model
outputs, with particular attention to the O3 construction described in
\textcite{kalisz2026spendoraclebudgetpractical} and its relationship to the
broader optimization setting in
\textcite{willis2026sampleefficientoptimisationoutputsgenerative}.

The scope is deliberately limited to three questions:
\begin{enumerate}
    \item Why is direct optimization in the original Boltz-2 latent space
          inefficient under a small oracle budget?
    \item How does O3 construct a lower-dimensional, non-redundant search
          space while preserving the intended latent statistics?
    \item What assumptions are required for Bayesian optimization to use
          that space reliably?
\end{enumerate}

The source expects an external \texttt{references.bib} containing the citation
keys used below.  Bibliographic metadata is intentionally not reconstructed in
this report.

\section{Problem formulation}
\label{sec:problem}

Let \(D\) denote the dimension of the original latent representation.  For a
fixed conditioning input \(c\), a frozen generative model maps a latent
variable to an output structure:
\begin{equation}
    g:\latent\times\mathcal{C}\longrightarrow\structure,
    \qquad x=g(z;c).
    \label{eq:generator}
\end{equation}

The latent variable is sampled from the model's intended distribution.  In the
Gaussian case considered here,
\begin{equation}
    z\sim\mathcal{N}(0,I_D),
    \qquad z\in\latent=\R^D.
    \label{eq:latent-distribution}
\end{equation}

An expensive oracle assigns a scalar score to the generated output:
\begin{equation}
    r:\structure\longrightarrow\R,
    \qquad y=r(x)=r\bigl(g(z;c)\bigr).
    \label{eq:oracle}
\end{equation}

The original black-box problem is therefore
\begin{equation}
    \max_{z\in\R^D} r\bigl(g(z;c)\bigr).
    \label{eq:original-objective}
\end{equation}

For the 1CLL protein experiment, \(c\) is the fixed protein-sequence
conditioning information and the oracle is the TM-score against a reference
structure \(x_{\mathrm{GT}}\):
\begin{equation}
    r(x)=\TMScore(x,x_{\mathrm{GT}}),
    \qquad
    \max_{z\in\R^D}\TMScore\bigl(g(z;c),x_{\mathrm{GT}}\bigr).
    \label{eq:tm-score-objective}
\end{equation}

\section{Why direct latent optimization is difficult}
\label{sec:direct-optimization}

In the 1CLL configuration, the paper-level latent dimension is
\begin{equation}
    D=3552.
    \label{eq:latent-dimension}
\end{equation}
An optimizer that works directly with \(z\in\R^{3552}\) must infer a useful
objective landscape from a very small number of oracle evaluations.  When the
budget is \(N=100\), a Gaussian-process surrogate has far fewer observations
than would normally be needed to resolve arbitrary variation in all latent
coordinates.

There is also a distributional constraint.  A generative model is trained and
calibrated for latents with particular statistics.  An optimization procedure
that produces vectors with the wrong norm or scale may leave the typical input
regime even when the vectors are mathematically valid.  The surrogate space
must therefore reduce dimension without turning the latent generator into an
uncontrolled extrapolation problem.

\section{Construction of the O3 surrogate latent space}
\label{sec:construction}

The construction uses \(d\) selected seed latents.  This report uses \(d\) for
the number of seeds throughout; it does not reuse \(K\) for that purpose.  The
repository configuration uses uppercase \(K\) for the number \(k\) of
returned candidates, which is a different quantity.

\subsection{Random low-dimensional embeddings}

A standard random-subspace method forms a full-dimensional latent from a
lower-dimensional weight vector:
\begin{equation}
    z^\star=Zw,
    \qquad Z\in\R^{D\times d},\quad w\in\R^d.
    \label{eq:random-embedding}
\end{equation}

If the columns of \(Z\) are independent standard Gaussian vectors, then
\begin{equation}
    Zw\sim\mathcal{N}\bigl(0,\norm{w}^2 I_D\bigr).
    \label{eq:embedding-scale}
\end{equation}
Thus, arbitrary weights alter the scale of the latent distribution.  In the
Gaussian setting, the Latent Optimal Linear combinations construction removes
that scale mismatch by normalizing the weights:
\begin{equation}
    z^\star=\frac{Zw}{\norm{w}},
    \qquad w\ne 0.
    \label{eq:lol-normalization}
\end{equation}

\subsection{Removing scaling redundancy}

The normalized representation is invariant to positive rescaling:
\begin{equation}
    \frac{Zw}{\norm{w}}
    =\frac{Z(\alpha w)}{\norm{\alpha w}},
    \qquad \alpha>0.
    \label{eq:scaling-redundancy}
\end{equation}
Consequently, many different parameter vectors describe the same latent.  O3
removes this flat direction by restricting the weights to the unit sphere:
\begin{equation}
    w\in\sphere^{d-1}
    =\left\{v\in\R^d:\norm{v}=1\right\}.
    \label{eq:unit-sphere}
\end{equation}

When \(Z\) has full column rank, which occurs almost surely for independent
Gaussian seeds when \(d\leq D\), distinct unit weight vectors produce
distinct latents.  The unit-norm restriction therefore removes the scaling
redundancy while retaining the desired Gaussian scale.

\subsection{Removing curvature with a surrogate chart}

The sphere is not an unconstrained Euclidean domain.  O3 restricts the search
to the positive orthant
\begin{equation}
    \sphere^{d-1}_{+}
    =\left\{w\in\R^d:w_i\geq 0\ \text{for all }i,\ \norm{w}=1\right\}
    \label{eq:positive-sphere}
\end{equation}
and maps Euclidean coordinates into that region:
\begin{equation}
    \phi_{\mathrm{KR}}: \mathcal{U}\longrightarrow\sphere^{d-1}_{+},
    \qquad \mathcal{U}=[0,1]^{d-1}.
    \label{eq:kr-chart}
\end{equation}

The chart is one-to-one and smooth on the interior of the cube, with a
continuous extension to boundary points.  Boundary points represent limiting
cases in which one or more weights vanish; the implementation also uses the
boundary vertices to represent individual seed latents.  This convention
avoids claiming a globally smooth bijection on the closed cube while matching
the optimizer's box constraints.

\section{Hyperspherical geometry and the KR chart}
\label{sec:kr-chart}

\subsection{Why direction is the useful latent coordinate}

For high-dimensional Gaussian samples, the norm concentrates near
\(\sqrt{D}\).  Much of the meaningful variation is therefore directional.
For \(z_i=Zw_i\) and \(z_j=Zw_j\), with unit-norm weights and independent
Gaussian columns of \(Z\),
\begin{equation}
    \mathbb{E}\!\left[z_i^{\mathsf T}z_j\right]
    =D\,w_i^{\mathsf T}w_j.
    \label{eq:expected-inner-product}
\end{equation}

Concentration of the inner products and norms gives the approximation
\begin{equation}
    \frac{z_i^{\mathsf T}z_j}{\norm{z_i}\,\norm{z_j}}
    \approx w_i^{\mathsf T}w_j.
    \label{eq:cosine-approximation}
\end{equation}
The chart is consequently chosen so that Euclidean distance in \(\mathcal{U}\)
has an approximately monotone relationship with weight-space similarity:
\begin{equation}
    v\!\left(\norm{u_i-u_j}\right)
    \approx w_i^{\mathsf T}w_j,
    \qquad w_i=\phi_{\mathrm{KR}}(u_i),
    \label{eq:chart-similarity}
\end{equation}
where \(v\) is decreasing.  Exact distance preservation is impossible when a
curved manifold is flattened into a bounded Euclidean domain, so this is a
regularity objective rather than an exact isometry.

\subsection{Inverse-Beta stick breaking}
\label{sec:stick-breaking}

The repository implementation uses the measure-preserving KR chart.  For
\(u=(u_1,\ldots,u_{d-1})\in[0,1]^{d-1}\), define
\begin{equation}
    b_j=I^{-1}_{u_j}\!\left(\frac12,\frac{d-j}{2}\right),
    \qquad j=1,\ldots,d-1,
    \label{eq:inverse-beta}
\end{equation}
where \(I^{-1}_{p}(a,b)\) is the inverse regularized Beta cumulative
distribution function.  Let \(\rho_1=1\) and recursively define
\begin{equation}
    \rho_{j+1}=\rho_j(1-b_j),
    \qquad
    w_j=\sqrt{\rho_j b_j},
    \quad j=1,\ldots,d-1,
    \qquad
    w_d=\sqrt{\rho_d}.
    \label{eq:kr-weights}
\end{equation}

Every weight is nonnegative and
\begin{equation}
    \sum_{j=1}^{d}w_j^2=1.
    \label{eq:unit-weight-sum}
\end{equation}
The correct constraint is therefore a squared-sum constraint, not the simplex
constraint \(\sum_j w_j=1\).  If \(u\) is uniformly distributed on the cube,
the squared weights follow a symmetric Dirichlet distribution with parameter
\(1/2\), which gives the uniform distribution on the positive orthant of the
unit sphere.  The distribution-preserving Gaussian conclusion applies to
independently sampled seed latents; selecting seeds by their oracle scores is
an additional data-dependent step whose empirical effect must be reported
separately.

\subsection{Matrix orientation used by the implementation}
\label{sec:matrix-orientation}

The mathematical notation in \eqref{eq:random-embedding} stores the seeds as
columns:
\begin{equation}
    Z=[z_1\ \cdots\ z_d]\in\R^{D\times d},
    \qquad z(u)=Z\,\phi_{\mathrm{KR}}(u).
    \label{eq:column-seeds}
\end{equation}

The implementation stores the same data row-wise as
\(
S=Z^{\mathsf T}\in\R^{d\times D}
\), so its matrix operation is
\begin{equation}
    z(u)=S^{\mathsf T}w(u),
    \qquad\text{implemented as}\qquad
    \texttt{weights @ seed\_latents}.
    \label{eq:row-seed-implementation}
\end{equation}
The resulting vector still has dimension \(D\).  This transpose is only a
storage convention; it does not define a different latent-space mapping.

\section{Seed selection and oracle-budget allocation}
\label{sec:budget}

Let \(N\) be the total oracle budget, \(M\) the number of initial random
evaluations, \(d\) the number of selected seed latents, and \(k\) the number of
final candidates returned.  The construction has the following stages.

\begin{enumerate}
    \item Sample \(M\) independent latents \(z_i\sim\mathcal{N}(0,I_D)\),
          generate \(x_i=g(z_i;c)\), and evaluate \(y_i=r(x_i)\).
    \item Select the \(d\) highest-scoring seed latents and form
          \(Z=[z_{(1)}\ \cdots\ z_{(d)}]\).
    \item Reuse the already measured seed scores at the \(d\) chart vertices,
          because each vertex maps to one basis weight vector.
    \item Evaluate two fresh uniformly sampled points in \(\mathcal{U}\) to
          initialize the Bayesian-optimization model.
    \item Spend the remaining budget on sequential acquisition points.
\end{enumerate}

In the repository runner, the seed evaluations are already included in the
phase of \(M\) calls.  The actual number of new oracle calls is therefore
\begin{equation}
    M+2+(N-M-2)=N,
    \qquad
    \text{sequential acquisitions}=N-M-2.
    \label{eq:budget-accounting}
\end{equation}
The configuration must satisfy \(2\leq d\leq M<N\) and \(N-M\geq 2\).

The paper also describes inversion as an alternative seed-selection strategy:
latents may be obtained from known outputs using an approximate inverse of the
generative model.  The present report focuses on filtering the initial random
samples by oracle score; inversion is outside the current experimental scope.

\section{Bayesian optimization and deterministic generation}
\label{sec:bo}

After selecting the seeds, O3 optimizes the bounded objective
\begin{equation}
    F(u)=r\!\left(g\!\left(Z\phi_{\mathrm{KR}}(u);c\right)\right),
    \qquad u\in[0,1]^{d-1}.
    \label{eq:surrogate-objective}
\end{equation}
The high-dimensional Boltz-2 latent is created only after the optimizer
proposes a low-dimensional point.

The complete optimization loop is therefore
\begin{equation}
    u
    \xrightarrow{\ \phi_{\mathrm{KR}}\ } w
    \xrightarrow{\ Z\ } z
    \xrightarrow{\ g(\,\cdot\,;c)\ } x
    \xrightarrow{\ r\ } y
    \xrightarrow{\ \text{fit GP and maximize }\LogEI\ }
    u_{\mathrm{next}}.
    \label{eq:canonical-pipeline}
\end{equation}

\subsection{Gaussian-process surrogate}

Given observations
\begin{equation}
    \mathcal{D}_t=\{(u_i,y_i)\}_{i=1}^{t},
    \label{eq:bo-data}
\end{equation}
the optimizer models \(F\) with a Gaussian process:
\begin{equation}
    F\sim\mathcal{GP}\bigl(m(u),\kappa_\theta(u,u')\bigr).
    \label{eq:gp}
\end{equation}
The repository implementation uses a constant mean, an ARD RBF covariance
kernel, and standardized outcomes.  The next point is selected by Log
Expected Improvement:
\begin{equation}
    u_{\mathrm{next}}
    =\underset{u\in[0,1]^{d-1}}{\operatorname{arg\,max}}\;\LogEI(u).
    \label{eq:logei}
\end{equation}
After evaluating \(F(u_{\mathrm{next}})\), the new observation is appended to
\(\mathcal{D}_t\) and the process repeats until the budget is exhausted.

\subsection{Deterministic evaluation assumption}

The O3 experiment uses deterministic probability-flow sampling so that a fixed
latent and fixed conditioning input produce a stable generated structure:
\begin{equation}
    u_1=u_2
    \Longrightarrow
    z(u_1)=z(u_2),\quad
    g(z(u_1);c)=g(z(u_2);c),\quad
    F(u_1)=F(u_2).
    \label{eq:deterministic-evaluation}
\end{equation}

This is an experimental protocol assumption, not a universal requirement for
Bayesian optimization.  If the generator or oracle is stochastic, the
observations instead follow a noisy model such as
\begin{equation}
    y=F(u)+\varepsilon,
    \qquad \varepsilon\sim\mathcal{N}(0,\sigma^2),
    \label{eq:noisy-objective}
\end{equation}
and repeated sampling or an explicitly calibrated noisy-BO model may be
needed.  Under a small budget, uncontrolled sampler randomness can obscure
the spatial variation that the GP is intended to learn, which is why the
current O3 protocol fixes the sampler behavior.

\section{Experimental protocol}
\label{sec:protocol}

This section records the implementation details that must accompany the
mathematical description.  It is intentionally a protocol section rather than
a results section.

\subsection{Dataset and target}

The current target is 1CLL, using the fixed protein sequence and reference
chain specified by the project configuration.  The sequence contains 144
residues.  \TODO{insert the final dataset/reference citation and any target
preprocessing details.}

\subsection{Model and latent representation}

The generator is the frozen Boltz-2 model conditioned on the target sequence.
The sampler input is configured with 1,184 coordinate slots, giving
\(1{,}184\times3=3{,}552\) latent values.  The configured sequence produces
1,134 real atoms; the remaining slots are masked padding and are removed before
the output structure is written and scored.  O3 uses deterministic PF-ODE
sampling, while any baseline comparison must state its sampler choice
explicitly.

\subsection{Configuration}

The currently configured sweep is shown below.  Here \(k\) denotes the number
of returned candidates; the configuration file names this field \(K\).

\begin{table}[ht]
    \centering
    \caption{Budget configurations.  Outcome columns are intentionally omitted
    until the complete experiment is available.}
    \label{tab:budgets}
    \begin{tabular}{@{}lrrrr@{}}
        \toprule
        Budget & \(N\) & \(k\) & \(M\) & \(d\) \\
        \midrule
        n20\_k2     & 20   & 2   & 10  & 6  \\
        n50\_k5     & 50   & 5   & 25  & 7  \\
        n100\_k10   & 100  & 10  & 50  & 5  \\
        n200\_k20   & 200  & 20  & 100 & 10 \\
        n500\_k50   & 500  & 50  & 250 & 30 \\
        n1000\_k100 & 1000 & 100 & 400 & 50 \\
        \bottomrule
    \end{tabular}
\end{table}

\TODO{add the final replicate count, random seeds, MSA policy, sampler steps,
and hardware/software versions used for the reported run.}

\subsection{Metrics and results placeholder}

The primary metric is TM-score against the 1CLL reference structure, with
higher values preferred.  The final report should compare O3 with the selected
baseline under matched budgets and replicate seeds.

\begin{table}[ht]
    \centering
    \caption{Results placeholder.  No numerical conclusions are drawn in this
    theory-focused version.}
    \label{tab:results-placeholder}
    \begin{tabular}{@{}lccc@{}}
        \toprule
        Budget & O3 best TM-score & Baseline best TM-score & Improvement \\
        \midrule
        n20\_k2     & \TODO{pending} & \TODO{pending} & \TODO{pending} \\
        n50\_k5     & \TODO{pending} & \TODO{pending} & \TODO{pending} \\
        n100\_k10   & \TODO{pending} & \TODO{pending} & \TODO{pending} \\
        n200\_k20   & \TODO{pending} & \TODO{pending} & \TODO{pending} \\
        n500\_k50   & \TODO{pending} & \TODO{pending} & \TODO{pending} \\
        n1000\_k100 & \TODO{pending} & \TODO{pending} & \TODO{pending} \\
        \bottomrule
    \end{tabular}
\end{table}

\section{Limitations and open questions}
\label{sec:limitations}

Several issues should remain explicit when interpreting the method:
\begin{enumerate}
    \item The positive-orthant chart covers a restricted region of the full
          sphere.  It improves parameterization and keeps weights
          nonnegative, but it does not explore every signed combination of
          seeds.
    \item Flattening a sphere into a cube necessarily distorts some distances.
          The KR chart is selected for useful distributional and local
          similarity properties, not exact global distance preservation.
    \item Selecting seeds by oracle score makes the seed set data-dependent.
          The ideal Gaussian calculation applies directly to independently
          sampled seeds; the effect of score-based selection needs empirical
          validation.
    \item Deterministic PF-ODE sampling makes the current objective easier to
          model, but conclusions may change under a stochastic sampler.
    \item The source excerpt supplied for this report does not contain the
          paper's complete experimental sections.  Claims about performance
          should therefore be added only after the full protocol and results
          are verified.
\end{enumerate}

\section{Conclusion}

O3 turns a high-dimensional latent optimization problem into a bounded,
low-dimensional search over \(u\in[0,1]^{d-1}\).  The KR chart produces
nonnegative unit-norm weights, those weights form a latent through
\(z=Zw\), and Bayesian optimization models the resulting objective with a
Gaussian process and LogEI.  The key implementation invariant is
\(\sum_iw_i^2=1\), not \(\sum_iw_i=1\).  The remaining question is empirical:
whether this structured search space improves oracle efficiency under the
completed experimental protocol.

\appendix

\section{Distribution-preserving derivation}
\label{app:distribution}

For completeness, define \(q_i=w_i^2\).  The stick-breaking construction in
\eqref{eq:kr-weights} produces nonnegative values satisfying
\(\sum_iq_i=1\).  With the inverse-Beta parameters in
\eqref{eq:inverse-beta}, the vector \(q\) follows
\begin{equation}
    q\sim\operatorname{Dirichlet}\!\left(\frac12,\ldots,\frac12\right).
    \label{eq:dirichlet-squared-weights}
\end{equation}
Taking componentwise square roots maps this distribution to the uniform
surface measure on \(\sphere^{d-1}_{+}\).  Hence \(w\) is nonnegative and has
unit Euclidean norm.

If the seed columns are independent standard Gaussian vectors, then for any
fixed unit-norm \(w\),
\begin{equation}
    Zw=\sum_{i=1}^{d}w_i z_i
    \sim\mathcal{N}\!\left(0,\left(\sum_{i=1}^{d}w_i^2\right)I_D\right)
    =\mathcal{N}(0,I_D).
    \label{eq:distribution-preservation}
\end{equation}
This calculation explains why the unit-sphere constraint preserves the
Gaussian scale.  It should not be read as proving that score-selected seeds
remain independent standard Gaussian samples; that data-dependent selection
effect belongs in the empirical validation.

\section{Dimension and budget checks}
\label{app:checks}

The dimensions used by the implementation are
\begin{equation}
    u\in\R^{d-1},
    \qquad w\in\R^d,
    \qquad Z\in\R^{D\times d},
    \qquad z=Zw\in\R^D.
    \label{eq:dimension-check}
\end{equation}
With row-wise storage \(S=Z^{\mathsf T}\), the equivalent operation is
\begin{equation}
    (1\times d)(d\times D)=1\times D,
    \qquad w^{\mathsf T}S=z^{\mathsf T}.
    \label{eq:row-dimension-check}
\end{equation}

For each configured run, the number of oracle calls is
\begin{equation}
    \underbrace{M}_{\text{random seed collection}}
    +\underbrace{2}_{\text{fresh BO initialization}}
    +\underbrace{(N-M-2)}_{\text{sequential acquisitions}}
    =N.
    \label{eq:appendix-budget-check}
\end{equation}

\printbibliography

\end{document}
